%% 06_noise.tex -- Chapter 6: Noise and the Noise--Power Trade-off
%% Status: OUTLINE (build stage 3). Replace the outline with the full chapter when it is written.
\chapter{Noise and the Noise--Power Trade-off}\label{ch:noise}
\chapterlead{Every microvolt of noise you remove costs current. The noise efficiency factor tells you how close a design is to the physical limit.}

\begin{specnote}[Chapter status: outline]
This chapter is an outline. It will be written in build stage 3 (Part II with Labs 3--5). The topics, the design question and the labs below are fixed; the derivations, worked examples and figures are still to come.
\end{specnote}

\begin{objectives}
\begin{itemize}
  \item compute thermal, flicker, shot and $kT/C$ noise and refer them to the amplifier input;
  \item integrate a noise spectrum over the signal band with the right noise bandwidth;
  \item derive the noise efficiency factor (NEF) and the power efficiency factor (PEF) and use them to compare designs;
  \item explain why a single bipolar or subthreshold device has $\NEF \approx 1$ and a differential pair about~2;
  \item set a current budget per channel from the \qty{1.5}{\uVrms} noise target.
\end{itemize}
\end{objectives}

\begin{insight}[Design question]
What is the minimum supply current that can meet the EEG-mode noise target, and how close to it can a real amplifier get?
\end{insight}

\section{Noise sources}\label{sec:ch6-1}
Thermal noise of resistors and channels, $1/f$ noise, shot noise of junction currents and $kT/C$ noise of sampled capacitors, with their spectral densities and the conditions where each dominates in a biopotential amplifier.

\begin{keyeq}[Noise efficiency factor]
\begin{equation}
  \NEF = V_\mathrm{ni,rms}\sqrt{\frac{2\,I_\mathrm{tot}}{\pi\, U_T\, 4kT\, \mathrm{BW}}}, \qquad \PEF = \NEF^2\,\Vdd
  \label{eq:nef}
\end{equation}
\end{keyeq}

\section{Input-referred noise and noise bandwidth}\label{sec:ch6-2}
Referring output noise to the input; the equivalent noise bandwidth of first- and second-order roll-offs ($\pi/2$ and $\pi/(2\sqrt2)$ factors); integrating $1/f$ noise between a sub-hertz high-pass and the low-pass corner.

\section{Noise efficiency factor and power efficiency factor}\label{sec:ch6-3}
$\NEF = V_\mathrm{ni,rms}\sqrt{2 I_\mathrm{tot}/(\pi U_T\, 4kT\, \mathrm{BW})}$, introduced by Steyaert and Sansen~\cite{steyaert1987} and made the standard yardstick by Harrison and Charles~\cite{harrison2003}; $\PEF = \NEF^2\Vdd$ adds the supply voltage. Published amplifiers plotted against both.

\section{The physical limit}\label{sec:ch6-4}
Why a single BJT or weak-inversion MOSFET gives $\NEF = 1$, a differential pair $\sqrt2$ to~2, and current-reuse topologies below~2. What the limit means for the \qty{10}{\micro\watt} per channel budget.

\begin{tryit}[Labs for this chapter]
Lab~4 (\cref{sec:lab4}) computes the NEF of your Harrison LNA versus bias current.
\end{tryit}

\paragraph{Read with this chapter} Steyaert and Sansen~\cite{steyaert1987}, Harrison and Charles~\cite{harrison2003}, Sarpeshkar~\cite{sarpeshkar}, Razavi~\cite{razavi} (Chapter on noise).
